% TOP_PROBLEM: {"id": 30007421, "title": "Height-averaged logarithmic bridge-length windows for critical honeycomb self-avoiding walks", "category_id": 19, "category": {"id": 19, "name": "probability", "display_name": "Probability", "description": "Problems involving probability theory, stochastic processes, and random structures.", "slug": "probability", "order_index": 19, "created_at": "2026-07-31T15:26:25.671Z"}, "status": "open"}
% FIRST_POSTED: 2026-10-06
% ATTEMPT_STATUS: solved
% OPENAI_CLAIM: {"schema_version":1,"source_repo":"https://github.com/openai/math","source_commit":"adc7f1241b42e322a6451854ab7e4b4c146bf78a","release_date":"2026-10-06","match":"full","resolution":"proved","independently_reviewed":false,"second_pass":"agreed","families":[{"family":"237","title":"The three-quarter exponent for honeycomb self-avoiding walk","subject":"Probability and statistical mechanics","manuscripts":[{"title":"Cylinder amplitudes and logarithmic bridge-length windows on the honeycomb lattice","pdf_path":"preprints/Cylinder-amplitudes-and-logarithmic-bridge-length-windows-on-the-honeycomb-lattice-September-26-2026/main.pdf","readme_path":"preprints/Cylinder-amplitudes-and-logarithmic-bridge-length-windows-on-the-honeycomb-lattice-September-26-2026/README.md","date":"2026-09-26","theorem_ref":"L:thm:finite (sections/introduction.tex:64); L:prop:box-moment (sections/10.tex:17); L:thm:window (sections/introduction.tex:34)","statement_sources":[{"path":"preprints/Cylinder-amplitudes-and-logarithmic-bridge-length-windows-on-the-honeycomb-lattice-September-26-2026/build/sections/introduction.tex","line":64,"label":"L:thm:finite"},{"path":"preprints/Cylinder-amplitudes-and-logarithmic-bridge-length-windows-on-the-honeycomb-lattice-September-26-2026/build/sections/10.tex","line":17,"label":"L:prop:box-moment"},{"path":"preprints/Cylinder-amplitudes-and-logarithmic-bridge-length-windows-on-the-honeycomb-lattice-September-26-2026/build/sections/introduction.tex","line":34,"label":"L:thm:window"}]}],"lean":{"doc_path":"lean/docs/237.md","comparators":[],"declarations":[],"covers_main_theorem":null},"reasoning_trace":null}]}
\documentclass[11pt]{article}
\usepackage[margin=25mm]{geometry}
\usepackage{fontspec}
\setmainfont{Latin Modern Roman}
\usepackage{amsmath,amssymb,mathtools}
\usepackage{unicode-math}
\setmathfont{Latin Modern Math}
\usepackage{enumitem}
\usepackage{xurl,hyperref,bookmark}
\hypersetup{colorlinks=true,urlcolor=blue}
\setcounter{secnumdepth}{0}
\setlength{\parindent}{0pt}
\setlength{\parskip}{5pt}
\setlength{\emergencystretch}{3em}
\begin{document}
\section*{Height-averaged logarithmic bridge-length windows for critical honeycomb self-avoiding walks}

\subsection{OpenAI's claim}
OpenAI presents a proof in the selected manuscript. After summing heights through order $h(\log h)^{1/64}$, a logarithmic window around $h^{4/3}$ is claimed to carry at least $h^{3/4}(\log h)^{-C}$ bridge mass.

\subsection{Mathematical statement of the claimed result}
\paragraph{Result 1.} At honeycomb critical fugacity \(X=(2+\sqrt2)^{-1/2}\), put \(\mu(\omega)=X^{|\omega|}\), counting visited vertices. For large real \(h\), define \(R_h=h(\log h)^{1/64}\), \(\ell_-=h^{4/3}(\log h)^{-1/8}\), \(\ell_+=h^{4/3}(\log h)^{1/2}\). The claim gives
\[
 \sum_{1\le H\le2R_h}\sum_{\omega\in\mathcal B_H,\ \ell_-\le|\omega|\le2\ell_+}\mu(\omega)
 \ge c h^{3/4}(\log h)^{-C}.
\]
This averages over integer bridge heights. Separate finite estimates are \(B_H\asymp H^{-1/4}\), arch mass \(A_g\asymp g^{-5/4}\), with a lower bound retained at diameter \(O(g)\). A fixed-width \(TH\) bridge corridor retains mass \(c_TH^{-1/4}\) and has length-weighted mass at most \(C_TH^{13/12}\), even after fixing either weak endpoint-displacement sign. For strip arches/bridges, diameter tail mass is at most \(CH^{-1/4}e^{-cr/H}\) for \(r\ge CH\). Polygon shells satisfy \(\sum_{R\le\operatorname{diam}P<2R}\mu(P)|P|^2\le CR^{2/3}(1+\log\log R)\). In convex finite triangle domains of diameter \(C_0r\) and bounded side count, all-endpoint chord length-weighted mass is at most \(Cr^{25/12}\). The polygon sum ranges over unrooted polygons modulo primitive lattice translations, so each translation class is counted once.

\subsection{Outline of the argument}
Finite-circumference loop expansions are recalibrated against positive cap amplitudes. Winding-pressure tails and divided-slice enclosing estimates transfer this normalization to strip crossings, concentrating a controlled mass in the specified length window after height averaging.

\subsection{Scope relative to this problem}
The logarithmic bridge-length window theorem is the literal estimate in the question, with R\_h = h(log h)\textasciicircum{}(1/64), lower length h\textasciicircum{}(4/3)(log h)\textasciicircum{}(-1/8), upper length 2h\textasciicircum{}(4/3)(log h)\textasciicircum{}(1/2), height range 1 <= H <= 2R\_h, and lower mass c h\textasciicircum{}(3/4)(log h)\textasciicircum{}(-C). It retains the fixed source, free terminal lateral port, center-count length, and height averaging. Inventory does not list this manuscript in Lean Scope 237; the family doc lists other companion papers, so its main Lean flag is null. Within these assumptions, this is a claimed affirmative answer to the stated question.

\subsection{Formal verification}
For family 237, the selected Lean statements give critical-strip mass and displacement exponents, bridge summability, and existence of pulled free energy. The $4/3$ bridge-length and $3/4$ spatial laws and the other polygon/chord exponents are outside these statements. No Comparator is assigned to the selected manuscript here; the family page selects other papers or supporting statements. This site has not executed the Lean code. The scope descriptions do not certify the remaining manuscript arguments.

\subsection{Status caveat}
OpenAI attributes these manuscripts to work by an internal model. Its release cautions that arguments without formalization can contain mistakes. This catalogue has not independently checked the proofs; the labels report the mathematical claims attributed to OpenAI.

\subsection{Sources}
\begin{itemize}
\item OpenAI, \emph{Cylinder amplitudes and logarithmic bridge-length windows on the honeycomb lattice}, 2026-09-26; citation key \texttt{OAI:Cylinder-amplitudes-and-logarithmic-bridge-length-windows-on-the-honeycomb-lattice-September-26-2026}. \href{https://github.com/openai/math/blob/adc7f1241b42e322a6451854ab7e4b4c146bf78a/preprints/Cylinder-amplitudes-and-logarithmic-bridge-length-windows-on-the-honeycomb-lattice-September-26-2026/main.pdf}{PDF (GitHub)}; \href{https://github.com/openai/math/raw/adc7f1241b42e322a6451854ab7e4b4c146bf78a/preprints/Cylinder-amplitudes-and-logarithmic-bridge-length-windows-on-the-honeycomb-lattice-September-26-2026/main.pdf}{PDF (direct)}; \href{https://github.com/openai/math/blob/adc7f1241b42e322a6451854ab7e4b4c146bf78a/preprints/Cylinder-amplitudes-and-logarithmic-bridge-length-windows-on-the-honeycomb-lattice-September-26-2026/README.md}{Manuscript page and citation}; \href{https://github.com/openai/math/blob/main/preprints/Cylinder-amplitudes-and-logarithmic-bridge-length-windows-on-the-honeycomb-lattice-September-26-2026/main.pdf}{Latest version}; Statement source: \href{https://github.com/openai/math/blob/adc7f1241b42e322a6451854ab7e4b4c146bf78a/preprints/Cylinder-amplitudes-and-logarithmic-bridge-length-windows-on-the-honeycomb-lattice-September-26-2026/build/sections/introduction.tex\#L64}{L:thm:finite}; Statement source: \href{https://github.com/openai/math/blob/adc7f1241b42e322a6451854ab7e4b4c146bf78a/preprints/Cylinder-amplitudes-and-logarithmic-bridge-length-windows-on-the-honeycomb-lattice-September-26-2026/build/sections/10.tex\#L17}{L:prop:box-moment}; Statement source: \href{https://github.com/openai/math/blob/adc7f1241b42e322a6451854ab7e4b4c146bf78a/preprints/Cylinder-amplitudes-and-logarithmic-bridge-length-windows-on-the-honeycomb-lattice-September-26-2026/build/sections/introduction.tex\#L34}{L:thm:window}.
\item Family 237: \href{https://github.com/openai/math/blob/adc7f1241b42e322a6451854ab7e4b4c146bf78a/lean/docs/237.md}{Lean scope}.
\item \href{https://github.com/openai/math/blob/adc7f1241b42e322a6451854ab7e4b4c146bf78a/CONTENTS.md}{OpenAI catalogue entry}.
\end{itemize}
\end{document}
